\chapter{Background and Positioning}
\label{ch:background}

This chapter fixes the four objects the rest of the thesis manipulates: the
Markov sequence that generates the data, the Ornstein--Uhlenbeck channel that
corrupts it, the time reversal that makes generation possible, and the identity
that turns the reversal's one unknown term into a posterior expectation. It
closes by stating what the surrounding literature settles and what it leaves
open.

The treatment is deliberately compressed. Every result below is standard, and
each is derived in full in \compendium{}: the variational calculus and the
statistical mechanics from which these models descend, the integrating-factor
solution of the Ornstein--Uhlenbeck equation, the Fokker--Planck route to the
reverse drift, and an extended review of the solvable-diffusion literature.
Nothing in the body depends on those derivations. The body depends only on the
statements collected here.

\section{Markov sequences}\label{sec:stationarity}\label{sec:markov}

A \emph{stochastic process} is a collection of random variables indexed by
time, $(X_k)_{k \geq 0}$ in discrete time or $(X_t)_{t \geq 0}$ in continuous
time, defined on a common probability space. Two notions of time homogeneity
are used constantly below and are worth separating precisely
\citep{brockwell1991time}.

\begin{definition}[Strict and weak stationarity]\label{def:stationarity}
A process $(X_k)$ is \emph{strictly stationary} if all its finite-dimensional
laws are invariant under time shifts: for every $n$, every set of indices
$k_1, \dots, k_n$, and every shift $h$,
\begin{equation*}
  (X_{k_1}, \dots, X_{k_n})
  \;\stackrel{d}{=}\;
  (X_{k_1 + h}, \dots, X_{k_n + h}) .
\end{equation*}
It is \emph{weakly (second-order) stationary} if its mean is constant and its
autocovariance depends only on the lag,
$\operatorname{Cov}(X_k, X_{k+d}) = \gamma(d)$.
\end{definition}

Strict stationarity implies weak stationarity whenever second moments exist,
and for \emph{Gaussian} processes the two coincide, because a Gaussian law is
determined by its first two moments. The distinction matters here because
Chapter~\ref{ch:laplace} replaces the Gaussian innovation with a Laplace one,
at which point the two notions separate and only the strict one is preserved by
the construction. Stationarity also has a matrix consequence the research
chapters exploit: the covariance of $L$ consecutive observations of a weakly
stationary process, $(\Sigma)_{ij} = \gamma(i - j)$, is constant along
diagonals, a symmetric \emph{Toeplitz} matrix. Wherever a covariance in this
thesis is Toeplitz, stationarity is the reason.

\begin{definition}[Markov chain]
A sequence of random variables $(X_k)_{k \geq 0}$ with values in a state space
$\mathcal{X}$ is a \emph{Markov chain} if for every $k$ and every measurable
$A \subseteq \mathcal{X}$,
\begin{equation}
  \Prob\big(X_{k+1} \in A \,\big|\, X_k, X_{k-1}, \dots, X_0\big)
  \;=\; \Prob\big(X_{k+1} \in A \,\big|\, X_k\big) .
  \label{eq:markov-property}
\end{equation}
The conditional law is encoded by a \emph{transition kernel}
$\kernel(x' \,|\, x)$: a density, or a matrix for discrete $\mathcal{X}$, with
$\int \kernel(x'|x)\,\dd x' = 1$ for every $x$.
\end{definition}

\begin{example}[The Gaussian AR(1) chain]\label{ex:ar1}
The running example of this thesis is the scalar Gaussian autoregression of
order one,
\begin{equation}
  a_{k+1} \;=\; \alpha\, a_k + \eta_k ,
  \qquad
  \eta_k \stackrel{\text{iid}}{\sim} \Nd(0, \ivar),
  \qquad k = 0, \dots, L-2 ,
  \label{eq:ar1}
\end{equation}
with transition kernel $\kernel(a' \mid a) = \Nd(a'; \alpha a, \ivar)$. Taking
$\ivar = 1 - \alpha^2$ and $a_0 \sim \Nd(0,1)$ makes it stationary with unit
marginal variance and autocovariance $\gamma(d) = \alpha^{|d|}$, so its
covariance is Toeplitz for the reason just given. Every chapter from
Chapter~\ref{ch:gaussian} onwards concerns \eqref{eq:ar1}, one of its
non-Gaussian variants, or what an estimator can recover of it through noise.
\end{example}

The Markov property \eqref{eq:markov-property} is exactly a statement about
\emph{factorisation}: the joint law of the first $L$ states is
\begin{equation}
  p(x_0, x_1, \dots, x_{L-1})
  \;=\; p_0(x_0)\, \prod_{k=0}^{L-2} \kernel(x_{k+1} \,|\, x_k) .
  \label{eq:chain-factorisation}
\end{equation}
This is the most heavily used equation of the thesis. It says the joint
distribution of a Markov trajectory is a product of \emph{local} terms, each
touching two consecutive variables: the structure that makes the chain a tree
in the factor-graph sense of Chapter~\ref{ch:graphs}, that makes its precision
matrix tridiagonal in the Gaussian case of Chapter~\ref{ch:gaussian}, and that
makes belief propagation exact. Marginals evolve by integrating the kernel,
$p_{k+1}(x') = \int \kernel(x'|x)\, p_k(x)\,\dd x$, and a law $\pi$ with
$\pi = \kernel \pi$ is \emph{stationary}: a chain started in $\pi$ remains in
it, and its trajectory is then strictly stationary in the sense of
Definition~\ref{def:stationarity}.

One change of perspective sets up the research problem. Markov chains are
classically studied as \emph{algorithms}, where the stationary law is the target
and the trajectory a means to reach it. From Chapter~\ref{ch:gaussian} onwards
the chain \eqref{eq:ar1} plays the opposite role: it is the
\emph{data-generating process}, one trajectory $(a_0, \dots, a_{L-1})$ is one
data point, and the object of interest is the law of the whole trajectory after
every coordinate has been independently corrupted. The factorisation
\eqref{eq:chain-factorisation} then stops being a computational convenience and
becomes the structure to be recovered.

\section{The Ornstein--Uhlenbeck corruption channel}\label{sec:ou}

The noising process used by diffusion models, and by every result of this
thesis, is the \emph{Ornstein--Uhlenbeck process}
\citep{uhlenbeck1930theory}: Brownian motion pulled back to the origin by a
linear spring,
\begin{equation}
  \dd X_t \;=\; -X_t\,\dd t \;+\; \sqrt{2}\;\dd W_t ,
  \qquad X_0 = a .
  \label{eq:ou-sde}
\end{equation}
The drift and diffusion constants are the variance-preserving normalisation of
the diffusion literature \citep{song2021scorebased}; any other choice rescales
time. The process is the unique Gaussian, Markov, stationary process in
continuous time, and it is exactly solvable. Solving \eqref{eq:ou-sde} by
integrating factor --- the derivation is routine and is carried out in
\compendium{} --- gives a Gaussian transition kernel,
$p(x_t \mid X_0 = a) = \Nd(x_t;\, a e^{-t},\, \Deltat)$ with
$\Deltat = 1 - e^{-2t}$. Corrupting a data point $a$ for a time $t$ therefore
produces
\begin{equation}
  x \;=\; a\, e^{-t} \;+\; \sqrt{\Deltat}\;\varepsilon,
  \qquad \varepsilon \sim \Nd(0, 1):
  \label{eq:ou-channel}
\end{equation}
the signal is \emph{contracted} by $e^{-t}$ while isotropic noise of variance
$\Deltat$ fills in. At $t = 0$ the channel is the identity; as $t \to \infty$,
$x \sim \Nd(0,1)$ regardless of $a$. The signal-to-noise ratio
$e^{-2t}/\Deltat$ sweeps continuously from $\infty$ to $0$ and diverges as
$1/2t$ for small $t$, a factor that reappears throughout the thesis as an error
amplifier. The channel interpolates between the data distribution and a
standard Gaussian, which is what a diffusion model needs.

Two structural remarks are used repeatedly later. First, the kernel is Gaussian
\emph{in $a$ as well as in $x$}: read as a function of the clean variable,
$\Nd(x; a e^{-t}, \Deltat) \propto
\exp[-(e^{-2t}/2\Deltat)(a - x e^{t})^2]$, a Gaussian likelihood factor of
precision $e^{-2t}/\Deltat$ centred at the deconvolved observation $e^t x$.
This is what makes the posteriors of Chapter~\ref{ch:gaussian} jointly Gaussian
and the belief-propagation messages closable in a finite-dimensional family.
Second, if $a \sim \Nd(\mu, \Sigma_0)$ is pushed through \eqref{eq:ou-channel}
coordinate-wise with \emph{independent} noises, the output is Gaussian with
\begin{equation}
  \mu_t = e^{-t}\mu,
  \qquad
  \Sigt \;=\; e^{-2t}\, \Sigma_0 \;+\; \Deltat\, I ,
  \label{eq:ou-cov-map}
\end{equation}
since means and covariances transform linearly and the added noise is white.
Equation \eqref{eq:ou-cov-map} --- shrink the covariance, add a multiple of the
identity --- is the entire forward process of this thesis in one line.

\section{Reverse diffusion}\label{sec:reversal}

Suppose data $X_0 \sim p_0$ are corrupted by a forward SDE
\begin{equation}
  \dd X_t \;=\; f(X_t, t)\,\dd t \;+\; g(t)\, \dd W_t
  \label{eq:generic-sde}
\end{equation}
up to time $T$, producing marginals $p_t$; the Ornstein--Uhlenbeck channel
\eqref{eq:ou-sde} is the case $f = -x$, $g = \sqrt{2}$. Can the corruption be
undone --- is there an SDE which, run from $X_T \sim p_T$ backwards, reproduces
the same marginals? The answer \citep{anderson1982reverse} is yes, and the
reverse drift needs exactly one object that the forward process does not
supply:
\begin{equation}
  \dd X_t \;=\;
  \Big[ f(X_t, t) \;-\; g(t)^2\, \nabla_x \log p_t(X_t) \Big]\,\dd t
  \;+\; g(t)\, \dd \bar W_t ,
  \label{eq:reverse-sde}
\end{equation}
where time runs backwards from $T$ to $0$ and $\bar W$ is a reverse-time Wiener
process. The correction to the drift is the \emph{score} $\nabla_x \log p_t$ of
the noised marginal. The forward process and the terminal law $p_T$ are known by
construction; the score is not, and estimating it is the entire modelling effort
of a diffusion model \citep{sohl2015deep, ho2020denoising, song2021scorebased}.
A Fokker--Planck argument making \eqref{eq:reverse-sde} plausible, together
with the probability-flow ODE it produces, is given in \compendium{}; the
rigorous statement is \citet{anderson1982reverse}. A recent book-length survey
of the model family is \citet{lai2025principles}.

\section{The score as a posterior expectation}\label{sec:tweedie}

Every result of this thesis passes through one rearrangement of the score.
Writing the noised marginal as the convolution of $p_0$ with the Gaussian
channel and differentiating under the integral
\citep{vincent2011connection}:
\begin{equation}
  \boxed{\;
  \nabla_x \log p_t(x)
  \;=\; \frac{e^{-t}\,\E[\,a \,|\, x\,] - x}{\Deltat}
  \quad\Longleftrightarrow\quad
  \E[\,a \,|\, x\,]
  = e^{t} \Big( x + \Deltat\, \nabla_x \log p_t(x) \Big).
  \;}
  \label{eq:tweedie}
\end{equation}
We refer to \eqref{eq:tweedie} throughout as the \emph{score--posterior
identity}; in the empirical-Bayes literature it is Tweedie's formula
\citep{robbins1956empirical, miyasawa1961empirical, efron2011tweedie}. The
score of the noisy marginal and the Bayes-optimal denoiser $\E[a|x]$ determine
each other exactly, with no assumption on $p_0$. Two readings recur. The score
points from $x$ towards the contracted image of the posterior mean, so denoising
direction and score direction coincide; and estimating the score \emph{is}
solving a Bayesian inference problem, not merely an analogue of one. For
sequence data with a Markov prior that inference problem lives on a tree, which
is why it is computable by message passing --- the programme made precise in
Chapter~\ref{ch:graphs} and executed in Chapter~\ref{ch:gaussian}.

For \emph{vector-valued} data $a \in \R^{L}$ noised coordinate-wise,
\eqref{eq:tweedie} holds componentwise with the \emph{full} posterior:
\begin{equation}
  \big( \nabla_x \log p_t(x) \big)_k
  \;=\; \frac{e^{-t}\, \E[\,a_k \,|\, x_0, \dots, x_{L-1}\,] - x_k}{\Deltat}.
  \label{eq:tweedie-joint}
\end{equation}
The conditioning is on \emph{all} coordinates. Even though the noise is
independent across frames, the posterior mean of frame $k$, and hence the $k$-th
score component, depends on every observation, because the prior couples the
frames. Equation \eqref{eq:tweedie-joint} is the precise reason the \emph{joint}
score differs from the per-frame marginal score
(Chapter~\ref{ch:development}), and Chapter~\ref{ch:gaussian} quantifies how
much coupling the conditioning induces as a function of $t$.

\section{Positioning}\label{sec:gap}\label{sec:related}

Five lines of work bound the problem this thesis addresses. Each is summarised
at the level of what it settles; the extended discussion, paper by paper, is in
\compendium{}.

\paragraph{Solvable diffusion models and phase transitions.} Treating the
reverse process as a stochastic dynamics on a high-dimensional energy landscape
has produced an exact picture for tractable priors: Gaussian mixtures and
random-energy targets exhibit sharp times at which the backward trajectory
commits to a class \citep{biroli2023generative, biroli2024dynamical,
raya2023spontaneous, achilli2026speciation}, and the sampling problem itself
admits a hardness analysis \citep{ghio2024sampling, cui2025solvable}. These
analyses fix the prior to a form with no internal sequential structure, so the
noised law factorises over coordinates or couples them only through a
low-dimensional latent.

\paragraph{Memorisation and generalisation.} A second strand asks when a
diffusion model reproduces its training set rather than the law behind it, and
locates the transition in the ratio of data to capacity
\citep{achilli2024losing, bonnaire2025memorize, kadkhodaie2024generalization,
favero2025bigger}. The mechanism is again characterised for priors without a
time axis internal to the data point.

\paragraph{Filtering and smoothing in state-space models.} Recovering a latent
Markov trajectory from noisy observations is classical: the Kalman filter
\citep{kalman1960new} and the Rauch--Tung--Striebel smoother
\citep{rauch1965maximum} solve it exactly in the linear-Gaussian case, and
expectation--maximisation \citep{dempster1977maximum} learns the transition when
it is unknown. This literature fixes the observation noise at a single level.
The diffusion setting instead requires the whole family indexed by $t$, and it
is the deformation across $t$, not the solution at any one $t$, that carries the
information a score model needs.

\paragraph{Exact inference on trees, and what happens off them.} Belief
propagation is exact on tree-shaped factor graphs \citep{pearl1988probabilistic,
mezard2009information}. On a chain with continuous variables the messages are
functions rather than vectors, and they close in a finite-dimensional family
only for special priors; moment projection \citep{minka2001} is the standard
response. What the projection costs, for a specific non-Gaussian chain, is not
generally known in closed form.

\paragraph{Networks as amortised inference.} Score networks are trained to
approximate a posterior expectation, which invites reading their architecture as
an inference procedure, and cascade-style analyses show structure being acquired
in stages \citep{bachtis2024cascade, sclocchi2025phase, favero2025compositional,
garnierbrun2026biased, zdeborova2016statistical}. The reading is suggestive
rather than established, and the comparison between a network and the exact
inference procedure for the same prior is rarely available, because the exact
procedure is rarely computable.

\medskip

Real sequential data --- video, audio, trajectories --- have exactly the
structure the solvable programme leaves out: an internal frame time with a
Markovian dependence along it. For such data the score \eqref{eq:tweedie-joint}
is a smoothing posterior on a chain, an object with seventy years of inference
theory behind it. Yet the questions a diffusion modeller asks about it are
answered neither by the classical literature, which fixes the noise level, nor
by the diffusion literature, which sets internal structure aside: how does the
coupling structure of $\nabla \log p_t$ deform with $t$; can it be computed
locally; what survives a Gaussian parametrisation of the messages when the data
are not Gaussian; can the transition itself be learned through the channel; and
how does the resulting estimator compare with a network trained on the same
data? Those five questions are the subject of
Chapters~\ref{ch:gaussian} to~\ref{ch:em-results}, with the inference tools
prepared in Chapter~\ref{ch:graphs} and the modelling choices that fixed the
formulation recorded in Chapter~\ref{ch:development}.
